
\section{Introduction}

\subsection{Motivation}

The parameter budget of a transformer language model is dominated by two 
components: the feed-forward network (FFN) and the self-attention mechanism. 
In DistilBERT, attention projections account for approximately \textbf{70\%} of the 
parameters in the transformer block. For GPT-2, this fraction rises to 
\textbf{75\%}. Reducing attention parameter cost is therefore a high-leverage 
compression strategy.

Standard multi-head self-attention (MHSA) stores four dense matrices per 
layer:

\[
W_Q, W_K, W_V, W_O \in \mathbb{R}^{d \times d}
\tag{1.1}
\]

totaling $4d^2$ parameters. For DistilBERT ($d = 768$), this is $4 \times 589{,}824 = 2{,}359{,}296$ 
parameters per layer, or 14.2M across 6 layers -- a substantial share of the 
66.9M total.

\textbf{Low-rank factorization} is a natural primitive for reducing this cost. 
For each $W$, one replaces it with a product:

\[
W \approx U V, \qquad
U \in \mathbb{R}^{d \times r}, \quad V \in \mathbb{R}^{r \times d}, \quad r \ll d
\tag{1.2}
\]

The parameter count drops from $d^2$ to $r(d + d) = 2rd$. The **reduction 
factor\textbf{ is $d / 2r$. For $d = 768$ and $r = 64$, this is }6$\times$**. Classical 
SVD (Eckart--Young, 1936) gives the optimal rank-$r$ approximation in 
Frobenius norm.

\textbf{The rank problem.} Classical SVD requires an a-priori choice of $r$, which 
is either:

\begin{itemize}
\item \textbf{Too small} -- the approximation error is high, and downstream accuracy
  drops.
\item \textbf{Too large} -- the parameter reduction is minimal, defeating the purpose.

\end{itemize}
In practice, different layers have different intrinsic ranks. Attention 
projections in early layers of DistilBERT tend to have rapidly-decaying 
spectra (effective rank ~30), while projections in middle layers decay 
slowly (effective rank ~100). A single global $r$ is suboptimal for both.

\subsection{Problem Statement}

Given a pretrained transformer $\mathcal{M}$ with $N_l$ layers and a target 
energy threshold $\eta \in (0, 1)$, we seek a compressed model $\mathcal{M}'$ 
such that:

\[
\mathcal{M}' \approx \mathcal{M}, \qquad
\frac{|\theta(\mathcal{M}')|}{|\theta(\mathcal{M})|} \ll 1
\tag{1.3}
\]

with the following constraints:

\begin{enumerate}
\item \textbf{Per-layer adaptive rank.} Different layers may have different ranks
  $r_\ell$ chosen automatically.
\item \textbf{Energy preservation.} For each factorized layer,
  $\sum_{i \leq r_\ell} \sigma_i^2 \geq \eta \sum_i \sigma_i^2$.
\item \textbf{Parameter reduction.} $r_\ell (d + d) < d^2$ for each layer.
\item \textbf{Multi-architecture support.} Must work on flat (DistilBERT), nested
  (RoBERTa), causal (GPT-2), and GQA (Qwen) transformers.
\item \textbf{No base retraining.} Only LoRA adapters may be trained.

\end{enumerate}
\subsection{Our Approach: MinimalNOVA}

\textbf{MinimalNOVA} is a drop-in replacement for the four attention projections 
of a transformer layer. It has three key components:

\begin{enumerate}
\item \textbf{SVD factorization.} Each $W$ is factorized as $U V$ via truncated SVD,
  with $U \in \mathbb{R}^{d_\text{out} \times r_\ell}$ and
  $V \in \mathbb{R}^{r_\ell \times d_\text{in}}$.

\item \textbf{Energy-adaptive rank.} Instead of fixing $r$, we choose:

  $$
  r\_\ell = \min\!\left(r\_\text{max}, \left|\left\{i : \sum\_{j \leq i} \sigma\_j^2 < \eta \sum\_j \sigma\_j^2\right\}\right| + 1\right)
  \tag{1.4}
  $$

  This selects the smallest rank that captures at least $\eta$ fraction of
  the total spectral energy, capped at $r_\text{max}$.

\item \textbf{LoRA recovery.} After factorization, the base factors $U, V$ are
  frozen and LoRA adapters $A, B$ are added:

  $$
  y = (x V^\top) U^\top + (x A^\top) B^\top \cdot s + b
  \tag{1.5}
  $$

  with $B = 0$ at initialization, ensuring the factorized output equals
  the SVD approximation at step 0.

\end{enumerate}
\subsection{Contributions}

We make the following contributions:

\begin{itemize}
\item \textbf{Energy-adaptive rank selection.} We propose a spectral-energy criterion
  that automatically selects a different rank per layer, avoiding the
  fixed-rank trade-off.
\item \textbf{A multi-architecture implementation.} We handle four attention families:
  flat (DistilBERT), nested (RoBERTa), causal with Conv1D (GPT-2), and GQA
  with RoPE (Qwen3.5). The nested-encoder bypass and Conv1D transposition
  are critical for correctness.
\item \textbf{A production-hardened pipeline.} Strict Frobenius validation, safe
  fallback to original layers, and detailed per-layer statistics.
\item \textbf{A systematic evaluation.} Benchmark on DistilBERT, RoBERTa, and GPT-2
  on SST-2 with ablations over rank cap $r_\text{max}$, energy threshold
  $\eta$, and LoRA rank.
\item \textbf{Full reproducibility.} Code, checkpoints, and scripts released under
  Apache 2.0.

\end{itemize}
\subsection{Paper Organization}

Section 2 reviews related work. Section 3 provides background on SVD, spectral 
energy, and LoRA. Section 4 presents MinimalNOVA. Section 5 describes 
integration with four architectures. Section 6 details the setup. Section 7 
reports results. Section 8 analyzes the energy-adaptive rank behavior. 
Section 9 discusses practical guidance. Section 10 lists limitations. 
Section 11 concludes.



\subsection{Notation}

\begin{lstlisting}

---
\end{lstlisting}